% Source PDF page 103; manual page 2-74.
\subsection{Relative Vehicle Calculations}\label{sec:relative-vehicle}

When more than one vehicle trajectory is computed, TAOS can express each vehicle's
position relative to another vehicle. For a selected current vehicle, the other
vehicle can be described either by relative range, azimuth, and elevation or by its
$x$, $y$, and $z$ coordinates in the current vehicle's body frame, as shown in
\cref{fig:relative-vehicle-calculations}.

\begin{figure}[htbp]
  \centering
  \resizebox{0.73\linewidth}{!}{\begin{tikzpicture}[>=Latex,line cap=round,line join=round,scale=0.9]
  \coordinate (o) at (-3.1,0);
  \coordinate (v) at (2.9,1.25);
  \node[left,align=center] at (o) {current\\vehicle};
  \fill (v) circle (2.2pt) node[right,align=left] {other\\vehicle};
  \draw[->,thick] (o)--(-2.0,1.15) node[above] {$\hat x_b$};
  \draw[->,thick] (o)--(2.55,-1.0) node[right] {$\hat y_b$};
  \draw[->,thick] (o)--(-3.1,-1.4) node[below] {$\hat z_b$};
  \draw[->,very thick] (o)--(v) node[pos=0.72,above] {$\Delta\vec r_i$};
  \draw[thin] (o)--(2.9,0.35)--(v);
  \draw[->] (-2.2,0.25) arc[start angle=18,end angle=42,radius=0.85]
    node[midway,above] {$\alpha_i$};
  \draw[->] (2.1,0.55) arc[start angle=0,end angle=48,radius=0.55]
    node[midway,right] {$\epsilon_i$};
  \draw (2.6,0.35)--(2.6,0.64)--(2.88,0.64);
\end{tikzpicture}
}
  \caption{Relative vehicle calculations.}
  \label{fig:relative-vehicle-calculations}
\end{figure}

The current vehicle's body-frame orientation is supplied by the guidance rules. It is
an input quantity and may change instantaneously, so TAOS does not compute relative
rates or accelerations in body coordinates.

The position of the $i$th vehicle relative to the current vehicle is
\begin{equation}
  \Delta\vec r_i=\vec r_i-\vec r,
  \label{eq:relative-vehicle-position}
\end{equation}
with associated unit vector
\begin{equation}
  \hat u_i=\frac{\Delta\vec r_i}{\lVert\Delta\vec r_i\rVert}.
  \label{eq:relative-vehicle-unit-vector}
\end{equation}
The relative azimuth and elevation are defined by
\begin{equation}
  \tan\alpha_i
  =\frac{\Delta\vec r_i\cdot\hat y_b}
         {\Delta\vec r_i\cdot\hat x_b},
  \label{eq:relative-vehicle-azimuth}
\end{equation}
\begin{equation}
  \sin\epsilon_i=-\hat u_i\cdot\hat z_b.
  \label{eq:relative-vehicle-elevation}
\end{equation}
The body-axis relative-position components are obtained directly by projecting
$\Delta\vec r_i$ onto $\hat x_b$, $\hat y_b$, and $\hat z_b$. The function
\texttt{vehicle\_obs} in the \texttt{path} module performs these calculations.
